\documentclass[../main]{subfiles}
\begin{document}

\chapter{Segundo Principio de la Termodinámica: Ciclo de Carnot}

\section{Segundo Principio de la Termodinámica}

El Segundo Principio de la Termodinámica establece que es imposible que un proceso cíclico, que no sea un ciclo de Carnot, tenga una eficiencia térmica mayor que la de un ciclo de Carnot entre los mismos focos térmicos\footnote{Focos térmicos son los cuerpos capaces de ceder o absorber cantidades de calor}, es decir, ningún motor térmico puede tener un rendimiento mayor que un motor reversible trabajando entre los mismos focos térmicos.

\subsection{Ciclo de Carnot}

El ciclo de Carnot es un modelo teórico de un motor térmico que opera entre dos fuentes de temperatura, una caliente y otra fría, a temperatura constante. Este ciclo consiste en cuatro procesos reversibles, dos isotérmicos y dos adiabáticos.

\info{Segundo Principio de la Termodinámica}{
El rendimiento de un motor térmico, es decir, la eficiencia con la que convierte el calor en trabajo, está dado por la expresión:
��
=
1
−
��
��
��
��
η=1− 
T 
c
​
 
T 
f
​
 
​
 
Donde:
\begin{itemize}
\item $\eta$: Eficiencia térmica del motor.
\item $T_f$: Temperatura de la fuente caliente.
\item $T_c$: Temperatura de la fuente fría.
\end{itemize}
}

\info{Ciclo de Carnot}{
El ciclo de Carnot se compone de cuatro procesos:
\begin{enumerate}
\item \textbf{Compresión isotérmica}: El gas se comprime isotérmicamente, absorbiendo calor de la fuente caliente.
\item \textbf{Expansión adiabática}: El gas se expande adiabáticamente, realizando trabajo sobre el sistema.
\item \textbf{Expansión isotérmica}: El gas se expande isotérmicamente, cediendo calor a la fuente fría.
\item \textbf{Compresión adiabática}: El gas se comprime adiabáticamente, realizando trabajo sobre el sistema.
\end{enumerate}
}

\info{Eficiencia del ciclo de Carnot}{
La eficiencia del ciclo de Carnot se calcula mediante la fórmula:
��
��
��
��
��
��
��
=
1
−
��
��
��
��
η 
Carnot
​
 =1− 
T 
c
​
 
T 
f
​
 
​
 
Donde:
\begin{itemize}
\item $\eta_{Carnot}$: Eficiencia del ciclo de Carnot.
\item $T_f$: Temperatura de la fuente caliente.
\item $T_c$: Temperatura de la fuente fría.
\end{itemize}
}

\actividad{Responder el siguiente cuestionario}
\begin{enumerate}
\item ¿Qué establece el Segundo Principio de la Termodinámica?
\item ¿Qué es un ciclo de Carnot y cuáles son sus componentes?
\item ¿Cuál es la expresión para calcular la eficiencia térmica de un motor?
\item ¿Cuál es la relación entre la eficiencia del ciclo de Carnot y las temperaturas de las fuentes térmicas?
% Agregar más preguntas
\end{enumerate}

\actividad{Resolver los siguientes problemas}
\begin{enumerate}
\item Calcular la eficiencia térmica de un motor que opera entre una fuente caliente a $500,K$ y una fuente fría a $300,K$.
\item Una máquina térmica opera entre dos fuentes a temperaturas de $600,K$ y $200,K$, respectivamente. Si la máquina produce $5000,J$ de trabajo, ¿cuánto calor se absorbe de la fuente caliente?
% Completar con los problemas faltantes
\end{enumerate}

\end{document}